\documentclass[10pt,a4paper]{article}
\usepackage[utf8]{inputenc}
\usepackage[T1]{fontenc}
\usepackage[ngerman]{babel}
\usepackage[margin=1.7cm,top=1.5cm,bottom=1.7cm]{geometry}
\usepackage{titlesec}
\titleformat*{\section}{\normalfont\large\bfseries}
\titleformat*{\subsection}{\normalfont\normalsize\bfseries}
\titlespacing*{\section}{0pt}{8pt}{3pt}
\titlespacing*{\subsection}{0pt}{6pt}{2pt}
\usepackage{amsmath,amssymb}
\usepackage{listings}
\usepackage{xcolor}
\usepackage{fancyhdr}
\usepackage{parskip}
\usepackage{needspace}

\newcommand{\mcbox}{\raisebox{-1pt}{\fbox{\rule{0pt}{7pt}\rule{7pt}{0pt}}}\hspace{6pt}}
\newcommand{\antwortlinien}[1]{%
  \par\vspace{2pt}%
  \newcount\zeilen \zeilen=#1
  \loop
    \noindent\rule{\linewidth}{0.4pt}\par\vspace{11pt}%
    \advance\zeilen by -1
  \ifnum\zeilen>0 \repeat
}
\lstset{language=Python,basicstyle=\ttfamily\footnotesize,keywordstyle=\bfseries,
  commentstyle=\itshape\color{gray},showstringspaces=false,frame=single,
  framerule=0.3pt,xleftmargin=8pt,framexleftmargin=6pt,aboveskip=6pt,belowskip=6pt}

\pagestyle{fancy}
\fancyhf{}
\fancyhead[L]{\small MDT5/2 -- Session 6: Autoencoder (AE, CAE, RandNet)}
\fancyhead[R]{\small Übungsaufgaben zur Klausurvorbereitung}
\fancyfoot[C]{\small Seite \thepage}
\renewcommand{\headrulewidth}{0.4pt}

\begin{document}

\begin{center}
  {\Large\bfseries Übungsaufgaben Session 6}\\[4pt]
  {\large Kapitel 7: Autoencoder, CAE, Ensembles (RandNet) -- Praktika 06 + 07}\\[4pt]
  {\footnotesize\itshape VAE entfällt (Klammer-Regel) -- nicht klausurrelevant.}\\[2pt]
  \small Stil der Übungssammlung MDT5/2 -- generierte Aufgaben, keine Originale
\end{center}

\vspace{4pt}

\section*{Aufgabe 1 -- Multiple Choice mit Begründung}
\emph{Markieren Sie jeweils die einzige richtige von drei Aussagen und begründen Sie Ihre Entscheidung.}

\subsection*{1.1 \ Für Autoencoder in der Novelty Detection gilt:}
\mcbox Der Anomalie-Score eines neuen Samples ist der Rekonstruktionsfehler; je höher,
desto wahrscheinlicher liegt eine Anomalie vor.

\mcbox Der Anomalie-Score ist der Wert des Codes $\boldsymbol{z}$ selbst.

\mcbox Autoencoder benötigen gelabelte Anomalien im Training.

Begründung:
\antwortlinien{2}

\subsection*{1.2 \ Für das Upsampling im Decoder eines Convolutional Autoencoders gilt:}
\mcbox \texttt{UpSampling2D} lernt seine Interpolationsgewichte im Training.

\mcbox Die transponierte Faltung (\texttt{Conv2DTranspose}) ermöglicht es, das Upsampling
mit lernbaren Filtern durchzuführen; mit Schrittweite 2 wird die Auflösung verdoppelt.

\mcbox Pooling-Schichten erhöhen die Auflösung im Decoder.

Begründung:
\antwortlinien{2}

\subsection*{1.3 \ Für Ensembles von Autoencodern gilt:}
\mcbox Alle Autoencoder des Ensembles müssen identisch aufgebaut und trainiert sein.

\mcbox Ensembles stabilisieren den Rekonstruktionsfehler, indem sich die Modelle z.\,B. in
Initialisierung, Trainingsdaten-Untermenge und Architektur unterscheiden.

\mcbox Der Ensemble-Score ist immer das Maximum aller Einzelfehler.

Begründung:
\antwortlinien{2}

\section*{Aufgabe 2 -- Convolutional Autoencoder in Keras schreiben}

Vervollständigen Sie den Code, indem Sie einen Convolutional Autoencoder anlegen.
Die Shape der Eingabedaten pro Sample beträgt $(32, 32, 3)$.
Der Encoder soll in \textbf{zwei} Faltungsschichten die Auflösung in Höhe und Breite
jeweils halbieren (\textbf{ohne} Pooling, also über die Schrittweite). Die Anzahl der Feature
Maps beträgt nach der ersten Schicht 8 und wird pro Schicht verdoppelt. Anschließend soll
ein Dense-Layer einen Code $\boldsymbol{z}$ mit 20 Elementen pro Sample erzeugen.
Der Decoder spiegelt den Encoder (transponierte Faltungen, Reihenfolge der
Feature-Map-Größen umgekehrt). Die Ausgabeschicht soll Werte zwischen 0 und 1 liefern.
Die Kernelgröße beträgt überall $3\times3$, als Aktivierung wird sonst immer ReLU verwendet.

\begin{lstlisting}
train_data = load_train_data()   # Werte bereits in [0; 1]

net =
\end{lstlisting}

\vspace{8.5cm}

\begin{lstlisting}
net.compile(optimizer='adam', loss='mse')
net.fit(x=train_data, y=train_data,
        batch_size=128, epochs=100, shuffle=True)
\end{lstlisting}

a) Warum wird bei \texttt{fit} zweimal \texttt{train\_data} übergeben (als \texttt{x} und
als \texttt{y})?
\antwortlinien{2}

b) Wie berechnet man nach dem Training den Anomalie-Score für neue Samples?
Beschreiben Sie die nötigen Schritte (kein Code erforderlich).
\antwortlinien{3}

\section*{Aufgabe 3 -- Identitätsfalle und Architekturwahl}

a) Warum besteht bei Autoencodern die Gefahr, dass die Identitätsfunktion gelernt wird,
und mit welcher Architekturentscheidung wird sie beim klassischen AE reduziert?
\antwortlinien{3}

b) Warum sind Convolutional Autoencoder für Bilder besser geeignet als reine
Dense-Autoencoder? Nennen Sie zwei Gründe.
\antwortlinien{2}

c) Nennen Sie eine traditionelle Anwendung des Autoencoders außerhalb der
Anomalieerkennung (aus der Vorlesung) und beschreiben Sie sie in einem Satz.
\antwortlinien{2}

\section*{Aufgabe 4 -- RandNet (Autoencoder-Ensemble)}

a) RandNet variiert drei Aspekte zufällig, damit die Ensemble-Mitglieder unterschiedliche
Information beitragen. Welche?
\antwortlinien{3}

b) Wie werden die zufällig \glqq gelöschten\grqq{} Verbindungen technisch umgesetzt
(Stichwort Maske)?
\antwortlinien{3}

c) Wie entsteht der finale Anomalie-Score eines neuen Samples aus den 100 Netzen?
(Fehlerberechnung, Normalisierung, Zusammenfassung)
\antwortlinien{3}

\vspace{10pt}
\begin{center}
\footnotesize\itshape
Nach Bearbeitung: Antworten an Claude schicken (Foto genügt) -- Korrektur mit Begründungs-Check.
\end{center}

\end{document}
